% =====================================================
\part{基礎－行列と連立一次方程式－}
% =====================================================

% =====================================================
\chapter*{はじめに ― なぜ線形代数を学ぶのか}
% =====================================================

大学に入って最初に出会う数学科目の代表が「微分積分学」と「線形代数学」である。
微分積分は高校数学の延長として姿が想像しやすいのに対し、線形代数は
「行列」「ベクトル空間」など聞き慣れない言葉が並び、何をしているのか
見失いやすい。そこで本書では、つねに次の一本の柱を意識して進む。

\begin{itembox}[l]{本書の柱}
線形代数とは、\textbf{連立一次方程式}と、それを生み出す
\textbf{「まっすぐな変換」(線形変換)}を統一的に扱う数学である。
\end{itembox}

連立一次方程式は中学以来の知り合いである。たとえば
\[
\begin{cases}
2x + y = 5 \\
x - 3y = -1
\end{cases}
\]
を解くことは誰でもできる。しかし、未知数が $100$ 個、式が $100$ 本あったら?
コンピュータに解かせるには、どんな原理で、どう手順化すればよいか?
解が「ちょうど一つ」「無数にある」「存在しない」の三択になるのはなぜか?
これらの問いに完全な答えを与えるのが線形代数であり、
その中心的な道具が本書の主役 \textbf{行列(matrix)} である。

線形代数は数学の中だけの話ではない。物理学(量子力学はほぼ線形代数で書かれる)、
統計学・データサイエンス、経済学、コンピュータグラフィックス、
そして近年の機械学習・AI は、すべて行列の計算を土台にしている。
大学で何を専攻するにせよ、線形代数は「読み書きそろばん」に相当する基礎体力である。

\bigskip
\noindent
【本書の前提】高校数学(数学I・A・II・B・III 程度)は理解しているものとする。
高校で「ベクトル」を学んだ経験は活きるが、行列の知識は一切仮定しない。

% =====================================================
\chapter{ベクトルを見直す ― 数の組という見方}
% =====================================================

\section{高校のベクトルと大学のベクトル}

高校では、ベクトルを「向きと大きさをもつ量」、すなわち矢印として学んだ。
大学の線形代数では、視点を少しずらして、ベクトルを
\textbf{数を縦に並べた組}とみなすことから出発する。

\begin{teigi}[数ベクトル]
実数を $n$ 個、縦に並べたもの
\[
\vect{a} = \mat{a_1 \\ a_2 \\ \vdots \\ a_n}
\qquad (a_1, a_2, \dots, a_n \in \R)
\]
を \textbf{$n$ 次元(数)ベクトル}という。$n$ 次元ベクトル全体の集合を
$\R^n$ と書く。
\end{teigi}

たとえば平面上の点 $(3, 2)$ に対応する矢印(位置ベクトル)は
\[
\vect{a} = \mat{3 \\ 2} \in \R^2
\]
と書ける。高校の矢印の世界は $\R^2$(平面)と $\R^3$(空間)であった。
大学では $n$ を $4, 5, \dots, 100, \dots$ と自由に大きくする。
「$100$ 次元の矢印」を目で見ることはできないが、
\textbf{数の組としてなら何次元でも同じように扱える}。
これが「矢印」から「数の組」へ視点を変える最大の利点である。

\begin{rei}
ある学生の期末試験の得点が、英語 $72$ 点、数学 $90$ 点、国語 $65$ 点、
理科 $88$ 点、社会 $70$ 点だったとする。この成績は
\[
\vect{x} = \mat{72 \\ 90 \\ 65 \\ 88 \\ 70} \in \R^5
\]
という $5$ 次元ベクトル一つで表せる。データを数の組=ベクトルとして
扱う発想は、統計学や機械学習の出発点でもある。
\end{rei}

\section{ベクトルの演算}

ベクトルの和とスカラー倍(実数倍)は、高校と同じく成分ごとに定める。

\begin{teigi}[和とスカラー倍]
$\vect{a} = \mat{a_1 \\ \vdots \\ a_n}$,
$\vect{b} = \mat{b_1 \\ \vdots \\ b_n}$, $c \in \R$ に対し、
\[
\vect{a} + \vect{b} = \mat{a_1 + b_1 \\ \vdots \\ a_n + b_n},
\qquad
c\,\vect{a} = \mat{c a_1 \\ \vdots \\ c a_n}.
\]
\end{teigi}

すべての成分が $0$ のベクトルを \textbf{零ベクトル}といい、$\vect{0}$ と書く。

実は、線形代数で起こることのほとんどは、この
「足す」と「実数倍する」のたった二つの演算の組合せ
\[
c_1 \vect{a}_1 + c_2 \vect{a}_2 + \cdots + c_k \vect{a}_k
\]
(これを $\vect{a}_1, \dots, \vect{a}_k$ の\textbf{線形結合}という)
だけから生まれる。「線形」代数という名前の由来はここにある。

\begin{mondai}\label{q:vec}     % 第1章の問
$\vect{a} = \mat{1 \\ 2}$, $\vect{b} = \mat{3 \\ -1}$ のとき、
$2\vect{a} - \vect{b}$ を計算せよ。また、
$\mat{5 \\ 0}$ を $\vect{a}$ と $\vect{b}$ の線形結合
$s\vect{a} + t\vect{b}$ の形で表せ。
\end{mondai}

後半は $s + 3t = 5$, $2s - t = 0$ という連立一次方程式を解くことに
ほかならない。\textbf{「線形結合で表せるか」という問いは連立一次方程式に帰着する}
――この事実は本書全体で繰り返し現れる。

% =====================================================
\chapter{行列の登場 ― 連立方程式の「係数」を取り出す}
% =====================================================

\section{行列の定義}

次の連立一次方程式を見てほしい。
\[
\begin{cases}
2x + y = 5 \\
x - 3y = -1
\end{cases}
\]
この方程式の「個性」を決めているのは、未知数の文字 $x, y$ ではなく、
係数 $2, 1, 1, -3$ と右辺 $5, -1$ である。
そこで係数だけを、位置関係を保ったまま長方形に並べてみる。
\[
\mat{2 & 1 \\ 1 & -3}
\]
このような「数の長方形配列」が行列である。

\begin{teigi}[行列]
$mn$ 個の数を、縦 $m$ 行・横 $n$ 列の長方形に並べたもの
\[
A = \mat{
a_{11} & a_{12} & \cdots & a_{1n} \\
a_{21} & a_{22} & \cdots & a_{2n} \\
\vdots & \vdots &        & \vdots \\
a_{m1} & a_{m2} & \cdots & a_{mn}
}
\]
を \textbf{$m \times n$ 行列}($m$ 行 $n$ 列の行列)という。
横の並びを\textbf{行}(row)、縦の並びを\textbf{列}(column)と呼び、
第 $i$ 行・第 $j$ 列にある数 $a_{ij}$ を $(i,j)$ \textbf{成分}という。
\end{teigi}

\begin{chui}
「行」と「列」はよく混同される。漢字の形で覚えるとよい:
「行」の字には横棒(二)が、「列」の字には縦棒(リ)が含まれる。
\textbf{行は横、列は縦}である。また、添字は必ず「行が先、列が後」
($a_{ij}$ は第 $i$ 行・第 $j$ 列)と約束する。
\end{chui}

\begin{rei}
\[
A = \mat{1 & 0 & -2 \\ 3 & 5 & 4}
\]
は $2 \times 3$ 行列で、$(1,3)$ 成分は $-2$、$(2,2)$ 成分は $5$ である。
\end{rei}

行と列の数が等しい $n \times n$ 行列を \textbf{$n$ 次正方行列}という。
また、$n$ 次元ベクトルは $n \times 1$ 行列とみなせる。
つまり\textbf{ベクトルは行列の特別な場合}である。

\section{行列の和とスカラー倍}

行列の和とスカラー倍は、ベクトルと同様に成分ごとに定める。
同じ型(同じ $m \times n$)の行列どうしでのみ和が定義される。

\begin{rei}
\[
\mat{1 & 2 \\ 3 & 4} + \mat{0 & -1 \\ 5 & 2} = \mat{1 & 1 \\ 8 & 6},
\qquad
3 \mat{1 & 2 \\ 3 & 4} = \mat{3 & 6 \\ 9 & 12}.
\]
\end{rei}

ここまでは「数の表の足し算」にすぎず、何も難しいことはない。
行列が本領を発揮するのは、次の\textbf{積}からである。

% =====================================================
\chapter{行列の積 ― なぜあの奇妙な定義なのか}
% =====================================================

\section{まず「行列 × ベクトル」から}

行列の積は、初学者が最初につまずく場所である。
「成分ごとに掛ければいいのに、なぜあんな複雑な定義を?」
と誰もが思う。答えは明快で、
\textbf{連立一次方程式を $A\vect{x} = \vect{b}$ と書けるようにするため}である。

もう一度、例の方程式を見る。
\[
\begin{cases}
2x + y = 5 \\
x - 3y = -1
\end{cases}
\]
係数行列 $A$、未知数ベクトル $\vect{x}$、右辺ベクトル $\vect{b}$ を
\[
A = \mat{2 & 1 \\ 1 & -3}, \qquad
\vect{x} = \mat{x \\ y}, \qquad
\vect{b} = \mat{5 \\ -1}
\]
とおく。このとき、方程式全体が
\[
A \vect{x} = \vect{b}
\]
という「一本の式」になるように積 $A\vect{x}$ を定義したい。
そのためには
\[
\mat{2 & 1 \\ 1 & -3} \mat{x \\ y} = \mat{2x + y \\ x - 3y}
\]
となっていなければならない。つまり、
\textbf{$A$ の各行とベクトル $\vect{x}$ を「内積のように」掛けて足す}
のが自然な定義である。

\begin{teigi}[行列とベクトルの積]
$m \times n$ 行列 $A = (a_{ij})$ と $n$ 次元ベクトル $\vect{x}$ に対し、
$A\vect{x}$ は次で定まる $m$ 次元ベクトルである:
\[
A\vect{x} = \mat{
a_{11} x_1 + a_{12} x_2 + \cdots + a_{1n} x_n \\
a_{21} x_1 + a_{22} x_2 + \cdots + a_{2n} x_n \\
\vdots \\
a_{m1} x_1 + a_{m2} x_2 + \cdots + a_{mn} x_n
}.
\]
\end{teigi}

これで、どんなに大きな連立一次方程式も
\[
A\vect{x} = \vect{b}
\]
という、中学校の一次方程式 $ax = b$ とそっくりな形に書ける。
「$a \neq 0$ なら $x = a^{-1}b$」という解法の類似物が行列の世界にもあるのか
――この問いが第 \ref{chap:inverse} 章の「逆行列」につながる。

\begin{mondai}\label{q:matvec}  % 第3章の問
$A = \mat{1 & 2 & 0 \\ -1 & 3 & 1}$,
$\vect{x} = \mat{2 \\ 1 \\ -3}$ のとき $A\vect{x}$ を計算せよ。
(答えは $2$ 次元ベクトルになるはずである。)
\end{mondai}

\section{行列を「関数」とみなす ― 線形変換}

積 $A\vect{x}$ の定義によって、行列 $A$ には新しい顔が生まれる。
$n$ 次元ベクトル $\vect{x}$ を入れると $m$ 次元ベクトル $A\vect{x}$ が
出てくる――これは $A$ が
\[
f(\vect{x}) = A\vect{x}
\]
という\textbf{関数(写像)}だということである。
高校で学んだ $f(x) = 2x$ のような関数の、入力と出力をベクトルにした版だと思えばよい。

この写像は次の二つの性質をもつ(積の定義から確かめられる):
\[
A(\vect{x} + \vect{y}) = A\vect{x} + A\vect{y},
\qquad
A(c\vect{x}) = c\,A\vect{x}.
\]
「和を保ち、定数倍を保つ」――この性質をもつ写像を\textbf{線形写像}
(特に $\R^n$ から $\R^n$ への場合は\textbf{線形変換})という。
直観的には、\textbf{原点を動かさず、直線を直線にうつす「まっすぐな変換」}である。

\begin{rei}[平面の回転]
平面のすべての点を、原点のまわりに角 $\theta$ だけ回転させる変換は、
行列
\[
R_\theta = \mat{\cos\theta & -\sin\theta \\ \sin\theta & \cos\theta}
\]
で表される。たとえば点 $\mat{1 \\ 0}$ を $90^\circ$ 回転すると
\[
\mat{\cos 90^\circ & -\sin 90^\circ \\ \sin 90^\circ & \cos 90^\circ}
\mat{1 \\ 0}
= \mat{0 & -1 \\ 1 & 0}\mat{1 \\ 0} = \mat{0 \\ 1}
\]
となり、確かに真上の点にうつる。
高校で天下り的に覚えた(あるいは加法定理で導いた)「回転の公式」は、
行列 $R_\theta$ を掛ける操作そのものだったのである。
\end{rei}

\begin{rei}[拡大・縮小と鏡映]
$\mat{2 & 0 \\ 0 & 2}$ はすべての点を原点から $2$ 倍に拡大する。
$\mat{1 & 0 \\ 0 & -1}$ は $x$ 軸に関する折り返し(鏡映)である。
各自、$\mat{1 \\ 1}$ などを掛けて確かめてみてほしい。
\end{rei}

\begin{itembox}[l]{行列の二つの顔}
\begin{itemize}
\item \textbf{静的な顔}:連立一次方程式の係数を並べた表(データ)
\item \textbf{動的な顔}:ベクトルをベクトルにうつす線形変換(操作)
\end{itemize}
同じ行列をこの二つの顔で自在に見られるようになることが、
線形代数習得の第一目標である。
\end{itembox}

\section{行列どうしの積 ― 「変換の合成」として}

行列 $A$ と $B$ の積 $AB$ はどう定義すべきか。
行列を変換と見る立場に立てば、答えは一つに決まる。
\textbf{「まず $B$ で変換し、続けて $A$ で変換する」という合成変換を
表す行列を $AB$ と呼ぶ}のである。すなわち、すべての $\vect{x}$ に対して
\[
(AB)\vect{x} = A(B\vect{x})
\]
が成り立つように定義する。この要請から成分の計算規則を導くと、
次の定義が得られる。

\begin{teigi}[行列の積]
$A$ を $m \times n$ 行列、$B$ を $n \times p$ 行列とする
(\textbf{$A$ の列数と $B$ の行数が一致}していることが必要)。
積 $AB$ は $m \times p$ 行列で、その $(i,j)$ 成分は
\[
(AB)_{ij} = a_{i1} b_{1j} + a_{i2} b_{2j} + \cdots + a_{in} b_{nj}
= \sum_{k=1}^{n} a_{ik} b_{kj}
\]
で与えられる。「$A$ の第 $i$ 行と $B$ の第 $j$ 列の内積」と覚えるとよい。
\end{teigi}

\begin{rei}
\[
\mat{1 & 2 \\ 3 & 4} \mat{0 & 1 \\ 5 & -1}
= \mat{1\cdot 0 + 2 \cdot 5 & 1 \cdot 1 + 2 \cdot (-1) \\
       3\cdot 0 + 4 \cdot 5 & 3 \cdot 1 + 4 \cdot (-1)}
= \mat{10 & -1 \\ 20 & -1}.
\]
\end{rei}

\begin{rei}[回転の合成と加法定理]
角 $\alpha$ の回転に続けて角 $\beta$ の回転を行うと、
合わせて角 $\alpha + \beta$ の回転になるはずである。
つまり $R_\beta R_\alpha = R_{\beta + \alpha}$ が成り立つはずで、
実際に左辺を積の定義で計算すると
\[
\mat{\cos\beta & -\sin\beta \\ \sin\beta & \cos\beta}
\mat{\cos\alpha & -\sin\alpha \\ \sin\alpha & \cos\alpha}
= \mat{
\cos\beta\cos\alpha - \sin\beta\sin\alpha &
-(\cos\beta\sin\alpha + \sin\beta\cos\alpha) \\
\sin\beta\cos\alpha + \cos\beta\sin\alpha &
\cos\beta\cos\alpha - \sin\beta\sin\alpha}
\]
となる。これが $R_{\alpha+\beta}$ に等しいという式は、
三角関数の\textbf{加法定理そのもの}である。
高校で苦労して証明した加法定理が、「回転の合成」という一言から
自然に出てくる――行列の威力を示す好例である。
\end{rei}

\section{積の性質 ― 数との違いに注意}

行列の積は、数の積と多くの性質を共有する。
$A, B, C$ の型が積を定義できるように整っているとき、
\[
(AB)C = A(BC), \qquad
A(B + C) = AB + AC, \qquad
(A + B)C = AC + BC.
\]
しかし、決定的に違う点が二つある。

\begin{itembox}[l]{数と行列の重要な違い}
\begin{enumerate}
\item \textbf{交換法則は成り立たない}:一般に $AB \neq BA$。
\item \textbf{零因子が存在する}:$A \neq O$, $B \neq O$ でも $AB = O$
となることがある($O$ はすべての成分が $0$ の\textbf{零行列})。
\end{enumerate}
\end{itembox}

\begin{rei}
$A = \mat{0 & 1 \\ 0 & 0}$, $B = \mat{1 & 0 \\ 0 & 0}$ とすると
\[
AB = \mat{0 & 0 \\ 0 & 0} = O, \qquad
BA = \mat{0 & 1 \\ 0 & 0} \neq O.
\]
$AB \neq BA$ であり、しかも $A \neq O$, $B \neq O$ なのに $AB = O$ である。
\end{rei}

このため、数の感覚で「$AB = O$ ならば $A = O$ または $B = O$」と
推論してはならないし、$(A+B)^2 = A^2 + 2AB + B^2$ も一般には成り立たない
(正しくは $(A+B)^2 = A^2 + AB + BA + B^2$)。
変換として考えれば交換不可能性は自然である:
「$90^\circ$ 回転してから $x$ 軸で折り返す」と
「$x$ 軸で折り返してから $90^\circ$ 回転する」では結果が異なる。

\section{単位行列}

数の世界の $1$(何を掛けても変えない数)に対応するのが単位行列である。

\begin{teigi}[単位行列]
対角成分(左上から右下への並び)がすべて $1$、他がすべて $0$ の
$n$ 次正方行列
\[
E = \mat{1 & 0 & \cdots & 0 \\ 0 & 1 & \cdots & 0 \\
\vdots & & \ddots & \vdots \\ 0 & 0 & \cdots & 1}
\]
を \textbf{単位行列}という($I$ と書く流儀もある)。
任意の $n$ 次正方行列 $A$ に対して $AE = EA = A$ が成り立つ。
\end{teigi}

変換として見れば、$E$ は「何もしない変換」(恒等変換)である。

% =====================================================
\chapter{連立一次方程式を解く ― 掃き出し法}\label{chap:gauss}
% =====================================================

\section{解法を「手順」にする}

中学・高校では連立方程式を加減法や代入法で、いわば臨機応変に解いてきた。
線形代数ではこれを、\textbf{誰がやっても(コンピュータがやっても)
同じ手順で必ず終わるアルゴリズム}に磨き上げる。
それが\textbf{掃き出し法}(ガウスの消去法)である。

使ってよい操作は次の三つだけである。
いずれも方程式の解を変えないことに注意してほしい。

\begin{itembox}[l]{行基本変形}
\begin{enumerate}
\item ある式を $c$ 倍する($c \neq 0$)
\item 二つの式を入れ替える
\item ある式に、別の式の $c$ 倍を加える
\end{enumerate}
\end{itembox}

計算の際は、文字 $x, y, z$ を毎回書くのは無駄なので、
係数と右辺だけを並べた\textbf{拡大係数行列}で作業する。

\begin{rei}[掃き出し法の実行]
\[
\begin{cases}
x + 2y + z = 8 \\
2x + 3y - z = 1 \\
-x + y + 2z = 7
\end{cases}
\]
を解く。拡大係数行列は
\[
\left(\begin{array}{ccc|c}
1 & 2 & 1 & 8 \\
2 & 3 & -1 & 1 \\
-1 & 1 & 2 & 7
\end{array}\right)
\]
である。第 $1$ 行を使って第 $1$ 列の他の成分を消す
(第 $2$ 行から第 $1$ 行の $2$ 倍を引き、第 $3$ 行に第 $1$ 行を加える):
\[
\left(\begin{array}{ccc|c}
1 & 2 & 1 & 8 \\
0 & -1 & -3 & -15 \\
0 & 3 & 3 & 15
\end{array}\right)
\]
第 $2$ 行を $-1$ 倍し、それを使って第 $2$ 列を掃き出す:
\[
\left(\begin{array}{ccc|c}
1 & 2 & 1 & 8 \\
0 & 1 & 3 & 15 \\
0 & 3 & 3 & 15
\end{array}\right)
\longrightarrow
\left(\begin{array}{ccc|c}
1 & 0 & -5 & -22 \\
0 & 1 & 3 & 15 \\
0 & 0 & -6 & -30
\end{array}\right)
\]
第 $3$ 行を $-\dfrac{1}{6}$ 倍し、第 $3$ 列を掃き出す:
\[
\left(\begin{array}{ccc|c}
1 & 0 & 0 & 3 \\
0 & 1 & 0 & 0 \\
0 & 0 & 1 & 5
\end{array}\right)
\]
よって解は $x = 3$, $y = 0$, $z = 5$ である。
\end{rei}

最後に得られた形(各行の先頭の $1$ が階段状に並び、
その上下が $0$ になっている形)を\textbf{簡約な階段形}という。
どんな行列も行基本変形で必ず簡約な階段形にでき、しかもその形は一通りに定まる。

\section{解は一つとは限らない}

掃き出し法のもう一つの偉さは、
\textbf{解が無数にある場合・存在しない場合も自動的に教えてくれる}ことである。

\begin{rei}[解が無数にある場合]
\[
\begin{cases}
x + y + z = 3 \\
2x + 2y + 2z = 6
\end{cases}
\]
第 $2$ 式は第 $1$ 式の $2$ 倍であり、実質的な条件は $1$ 本しかない。
掃き出すと第 $2$ 行はすべて $0$ になり、
$y = s$, $z = t$(任意の実数)とおけば解は
\[
\mat{x \\ y \\ z} = \mat{3 - s - t \\ s \\ t}
\qquad (s, t \text{ は任意の実数})
\]
と、$2$ 個のパラメータをもつ無数の解になる。
\end{rei}

\begin{rei}[解が存在しない場合]
\[
\begin{cases}
x + y = 1 \\
x + y = 2
\end{cases}
\]
を掃き出すと、$0 = 1$ という矛盾した行(左辺がすべて $0$ で右辺が $0$ でない行)が現れる。
このとき解は存在しない。幾何的には、平行な $2$ 直線が交点をもたないことに対応する。
\end{rei}

「ちょうど一つ」「無数」「なし」――解のあり方がこの三択に限られ、
どれになるかが掃き出し法で完全に判定できる。
そして「実質的な条件が何本あるか」を数える概念が、
後に学ぶ\textbf{階数(ランク)}である。

% =====================================================
\chapter{逆行列 ― 「割り算」を求めて}\label{chap:inverse}
% =====================================================

\section{逆行列の定義}

一次方程式 $ax = b$ は、$a \neq 0$ なら逆数 $a^{-1}$ を掛けて
$x = a^{-1} b$ と解ける。行列の方程式 $A\vect{x} = \vect{b}$ でも
同じことがしたい。「逆数」にあたるものを定義しよう。

\begin{teigi}[逆行列]
$n$ 次正方行列 $A$ に対して、
\[
AB = BA = E
\]
を満たす $n$ 次正方行列 $B$ が存在するとき、$A$ は\textbf{正則}であるといい、
$B$ を $A$ の\textbf{逆行列}と呼んで $A^{-1}$ と書く。
\end{teigi}

$A$ が正則なら、$A\vect{x} = \vect{b}$ の両辺に左から $A^{-1}$ を掛けて
\[
\vect{x} = A^{-1}\vect{b}
\]
と、解が一発で(しかもただ一つに)求まる。
変換の言葉でいえば、$A^{-1}$ は「$A$ の操作を元に戻す変換」である。
回転 $R_\theta$ の逆行列が逆回転 $R_{-\theta}$ であることは直観的にも明らかだろう。

ただし、数の $0$ に逆数がないのと同様、
\textbf{逆行列をもたない正方行列が存在する}。しかも行列の場合、
$O$ 以外にも正則でない行列がたくさんある点が数と異なる。

\begin{rei}
$A = \mat{1 & 2 \\ 2 & 4}$ は正則でない。実際、この行列の表す変換は
平面全体を直線 $y = 2x$ の上につぶしてしまう
(どんな $\vect{x}$ に対しても $A\vect{x}$ は $\mat{1 \\ 2}$ の定数倍になることを確かめよ)。
一度つぶれた情報は復元できないから、「元に戻す変換」は存在しえない。
\end{rei}

\section{$2$ 次正方行列の逆行列}

$2$ 次の場合には、逆行列の便利な公式がある。

\begin{teiri}[$2$ 次の逆行列公式]
$A = \mat{a & b \\ c & d}$ について、$ad - bc \neq 0$ のとき $A$ は正則で、
\[
A^{-1} = \frac{1}{ad - bc} \mat{d & -b \\ -c & a}.
\]
$ad - bc = 0$ のとき $A$ は正則でない。
\end{teiri}

正則かどうかの運命を握る量 $ad - bc$ を $A$ の\textbf{行列式}といい、
$\det A$ と書く。行列式は次章の主役である。

\begin{mondai}\label{q:inv2}    % 第5章の問
$A = \mat{2 & 1 \\ 1 & -3}$ の逆行列を求め、
冒頭の連立方程式 $A\vect{x} = \mat{5 \\ -1}$ を
$\vect{x} = A^{-1}\vect{b}$ により解け。
(答:$x = 2$, $y = 1$ となるはずである。)
\end{mondai}

\section{一般の逆行列の求め方}

$3$ 次以上では公式は複雑になるため、実際の計算には掃き出し法を流用する。
$A$ の右に単位行列 $E$ を並べた行列 $(A \mid E)$ に行基本変形を施し、
左半分が $E$ になったとき、右半分に現れているのが $A^{-1}$ である:
\[
(A \mid E) \longrightarrow (E \mid A^{-1}).
\]
途中で左半分に「すべて $0$ の行」が現れたら、$A$ は正則でない。

\begin{rei}
$A = \mat{1 & 2 \\ 3 & 4}$ について、
\[
\left(\begin{array}{cc|cc}
1 & 2 & 1 & 0 \\
3 & 4 & 0 & 1
\end{array}\right)
\longrightarrow
\left(\begin{array}{cc|cc}
1 & 2 & 1 & 0 \\
0 & -2 & -3 & 1
\end{array}\right)
\longrightarrow
\left(\begin{array}{cc|cc}
1 & 0 & -2 & 1 \\
0 & 1 & 3/2 & -1/2
\end{array}\right)
\]
よって $A^{-1} = \mat{-2 & 1 \\ 3/2 & -1/2}$ である。
(公式 $\det A = 1 \cdot 4 - 2 \cdot 3 = -2$ からも確かめられる。)
\end{rei}

% =====================================================
\chapter{行列式 ― 変換の「拡大率」}
% =====================================================

\section{行列式の幾何的な意味}

前章で現れた量 $\det A = ad - bc$ には、美しい幾何的意味がある。

\begin{itembox}[l]{行列式の意味}
$2$ 次正方行列 $A$ の表す線形変換によって、平面上の図形の面積は
$|\det A|$ 倍になる。すなわち\textbf{行列式は変換の面積拡大率}である
($3$ 次なら体積拡大率)。$\det A$ の符号は、変換が図形の「裏返し」を
含むかどうか(向きの反転)を表す。
\end{itembox}

実際、基本ベクトル $\vect{e}_1 = \mat{1 \\ 0}$, $\vect{e}_2 = \mat{0 \\ 1}$
の張る単位正方形(面積 $1$)は、変換 $A = \mat{a & b \\ c & d}$ によって
$\mat{a \\ c}$ と $\mat{b \\ d}$ の張る平行四辺形にうつる。
高校で学んだ平行四辺形の面積公式により、その面積は $|ad - bc|$ である。

この見方をすれば、「$\det A = 0$ ならば $A$ は正則でない」ことも腑に落ちる。
面積拡大率が $0$ とは、平面が直線(または $1$ 点)につぶされるということであり、
つぶす変換は元に戻せないからである。

\begin{teiri}[正則性の判定]
$n$ 次正方行列 $A$ について、次は同値である。
\begin{enumerate}
\item $A$ は正則である($A^{-1}$ が存在する)。
\item $\det A \neq 0$。
\item 連立方程式 $A\vect{x} = \vect{b}$ は任意の $\vect{b}$ に対して
ただ一つの解をもつ。
\end{enumerate}
\end{teiri}

この定理は、本書で学んだ三つの主役――逆行列・行列式・連立方程式――が
一つに結ばれる場所であり、線形代数前半のハイライトである
(同値条件はこの後の学習でさらに増えていく:階数、固有値、線形独立性…)。

\section{$3$ 次の行列式}

$3$ 次正方行列の行列式は次で定義される(サラスの方法):
\[
\det \mat{a_{11} & a_{12} & a_{13} \\ a_{21} & a_{22} & a_{23} \\ a_{31} & a_{32} & a_{33}}
= a_{11}a_{22}a_{33} + a_{12}a_{23}a_{31} + a_{13}a_{21}a_{32}
- a_{13}a_{22}a_{31} - a_{11}a_{23}a_{32} - a_{12}a_{21}a_{33}.
\]
「左上から右下への $3$ 本の対角線の積を足し、
右上から左下への $3$ 本の対角線の積を引く」と覚える。
ただし\textbf{サラスの方法は $3$ 次まで}しか使えない。
$4$ 次以上の行列式の定義と計算法(余因子展開・行基本変形の利用)は、
第II部で扱う。

\begin{mondai}\label{q:det3}    % 第6章の問
$\det \mat{1 & 2 & 1 \\ 2 & 3 & -1 \\ -1 & 1 & 2}$ を計算せよ。
これが $0$ でないことは、第 \ref{chap:gauss} 章の例の連立方程式の解が
ただ一つに定まったことと符合する。
\end{mondai}
